\documentclass[10pt,a4paper]{amsart}
\usepackage{amsmath,amssymb,amsthm,graphicx}
% Retained minimal preamble for the Stacks Project chapter more-morphisms.tex (More on Morphisms).
% Verbatim excerpt of the chapter preamble: Xy-pic setup (native preamble lines 14-18) and the native
% theorem declarations with the mathematical macros (native preamble lines 154-217). Copied mechanically;
% nothing added, removed or redefined. The wrapper class and the amsmath/amssymb/amsthm/graphicx base
% packages are supplied by the assembler, so the class/externaldocument part of the native preamble is omitted.
% For commutative diagrams we use Xy-pic
\usepackage[all]{xy}

% We use 2cell for 2-commutative diagrams.
\xyoption{2cell}
% Theorem environments.
%
\theoremstyle{plain}
\newtheorem{theorem}[subsection]{Theorem}
\newtheorem{proposition}[subsection]{Proposition}
\newtheorem{lemma}[subsection]{Lemma}

\theoremstyle{definition}
\newtheorem{definition}[subsection]{Definition}
\newtheorem{example}[subsection]{Example}
\newtheorem{exercise}[subsection]{Exercise}
\newtheorem{situation}[subsection]{Situation}

\theoremstyle{remark}
\newtheorem{remark}[subsection]{Remark}
\newtheorem{remarks}[subsection]{Remarks}

\numberwithin{equation}{subsection}

% Macros
%
\def\lim{\mathop{\mathrm{lim}}\nolimits}
\def\colim{\mathop{\mathrm{colim}}\nolimits}
\def\Spec{\mathop{\mathrm{Spec}}}
\def\Hom{\mathop{\mathrm{Hom}}\nolimits}
\def\Ext{\mathop{\mathrm{Ext}}\nolimits}
\def\SheafHom{\mathop{\mathcal{H}\!\mathit{om}}\nolimits}
\def\SheafExt{\mathop{\mathcal{E}\!\mathit{xt}}\nolimits}
\def\Sch{\mathit{Sch}}
\def\Mor{\mathop{\mathrm{Mor}}\nolimits}
\def\Ob{\mathop{\mathrm{Ob}}\nolimits}
\def\Sh{\mathop{\mathit{Sh}}\nolimits}
\def\NL{\mathop{N\!L}\nolimits}
\def\CH{\mathop{\mathrm{CH}}\nolimits}
\def\proetale{{pro\text{-}\acute{e}tale}}
\def\etale{{\acute{e}tale}}
\def\QCoh{\mathit{QCoh}}
\def\Ker{\mathop{\mathrm{Ker}}}
\def\Im{\mathop{\mathrm{Im}}}
\def\Coker{\mathop{\mathrm{Coker}}}
\def\Coim{\mathop{\mathrm{Coim}}}

% Boxtimes
%
\DeclareMathSymbol{\boxtimes}{\mathbin}{AMSa}{"02}

%
% Macros for moduli stacks/spaces
%
\def\QCohstack{\mathcal{QC}\!\mathit{oh}}
\def\Cohstack{\mathcal{C}\!\mathit{oh}}
\def\Spacesstack{\mathcal{S}\!\mathit{paces}}
\def\Quotfunctor{\mathrm{Quot}}
\def\Hilbfunctor{\mathrm{Hilb}}
\def\Curvesstack{\mathcal{C}\!\mathit{urves}}
\def\Polarizedstack{\mathcal{P}\!\mathit{olarized}}
\def\Complexesstack{\mathcal{C}\!\mathit{omplexes}}
% \Pic is the operator that assigns to X its picard group, usage \Pic(X)
% \Picardstack_{X/B} denotes the Picard stack of X over B
% \Picardfunctor_{X/B} denotes the Picard functor of X over B
\def\Pic{\mathop{\mathrm{Pic}}\nolimits}
\def\Picardstack{\mathcal{P}\!\mathit{ic}}
\def\Picardfunctor{\mathrm{Pic}}
\def\Deformationcategory{\mathcal{D}\!\mathit{ef}}
\usepackage{cleveref}
\crefname{originaldefinition}{original definition}{original definitions}
\Crefname{originaldefinition}{Original definition}{Original definitions}
\makeatletter
\newlabel{definition-weighting}{{0F3A}{source}{}{sourcepointer-1}{}}
\newlabel{definition-weighting@cref}{{[originaldefinition][0][]0F3A}{[1][0][]source}{}{}{}}
\makeatother
\crefname{originallemma}{original lemma}{original lemmas}
\Crefname{originallemma}{Original lemma}{Original lemmas}
\makeatletter
\newlabel{lemma-etale-makes-quasi-finite-finite-at-point}{{02LK}{source}{}{sourcepointer-2}{}}
\newlabel{lemma-etale-makes-quasi-finite-finite-at-point@cref}{{[originallemma][0][]02LK}{[1][0][]source}{}{}{}}
\makeatother
\crefname{originallemma}{original lemma}{original lemmas}
\Crefname{originallemma}{Original lemma}{Original lemmas}
\makeatletter
\newlabel{lemma-jumps-w}{{0F3G}{source}{}{sourcepointer-3}{}}
\newlabel{lemma-jumps-w@cref}{{[originallemma][0][]0F3G}{[1][0][]source}{}{}{}}
\makeatother
\crefname{originallemma}{original lemma}{original lemmas}
\Crefname{originallemma}{Original lemma}{Original lemmas}
\makeatletter
\newlabel{lemma-weighting-specialization}{{lemma-weighting-specialization}{source}{}{sourcepointer-4}{}}
\newlabel{lemma-weighting-specialization@cref}{{[originallemma][0][]lemma-weighting-specialization}{[1][0][]source}{}{}{}}
\makeatother
\crefname{originalreference}{original reference}{original references}
\Crefname{originalreference}{Original reference}{Original references}
\makeatletter
\newlabel{morphisms-lemma-image-nowhere-dense-finite}{{03HX}{source}{}{sourcepointer-5}{}}
\newlabel{morphisms-lemma-image-nowhere-dense-finite@cref}{{[originalreference][0][]03HX}{[1][0][]source}{}{}{}}
\makeatother
\crefname{originalreference}{original reference}{original references}
\Crefname{originalreference}{Original reference}{Original references}
\makeatletter
\newlabel{topology-lemma-dense-in-constructible}{{topology-lemma-dense-in-constructible}{source}{}{sourcepointer-6}{}}
\newlabel{topology-lemma-dense-in-constructible@cref}{{[originalreference][0][]topology-lemma-dense-in-constructible}{[1][0][]source}{}{}{}}
\makeatother
\title{A weighting vanishing on a dense subset vanishes identically}
\author{The Stacks Project Authors}
\date{stacks-project:more-morphisms.tex:local-snapshot-20261011; GNU FDL 1.2 or later}
\begin{document}
\maketitle
\noindent\textit{Editorial note.} Counted span is the author's own proof environment, source lines 22585-22616, immediately following the statement at 22577-22583.\par\medskip
\noindent\textit{Editorial note.} The proof opens with the author's suggestion 'We suggest reading Lemma [the specialisation lemma] first'; that sentence is retained verbatim.\par\medskip
\noindent\textit{Editorial note.} Two cross-chapter pointers (Morphisms, Topology) are recorded in assembly.external\_references.\par\medskip
\noindent\textit{Allowed standard prerequisites (editorial).}\begin{itemize}
\item Level sets of a weighting are locally constructible. chapter lemma lemma-jumps-w, cited at proof line 22604. Makes the zero locus of the pulled-back weighting a constructible set, so that density upgrades to containing a dense open.
\item The specialisation lemma for weightings. chapter lemma lemma-weighting-specialization, suggested at proof line 22586 and used through the argument. Lets vanishing be transported along specialisations, which is the mechanism behind the argument.
\item A constructible dense subset contains a dense open. Topology, lemma topology-lemma-dense-in-constructible, cited at proof line 22606. Turns density plus constructibility of the zero locus into density of an open piece.
\item Image of a nowhere dense set under a finite morphism. Morphisms, lemma morphisms-lemma-image-nowhere-dense-finite, cited at proof lines 22610-22611. Pushes the nowhere dense closed set Z down to a nowhere dense set in the base, so the integral vanishes on a dense open.
\end{itemize}\medskip
\ifdefined\hypertarget\hypertarget{sourcepointer-1}{}\fi
\noindent\textit{Source reference.} Original definition 0F3A: sources/stacks-more-morphisms/pinned-original-source.tex, source lines 22152--22152. Navigation pointer only: the counted proof cites this statement by \textbackslash{}ref as a standing input; the deduction the counted proof has to make is carried out inside the counted span, and no part of that deduction is delegated to this pointer.\par\medskip
\ifdefined\hypertarget\hypertarget{sourcepointer-2}{}\fi
\noindent\textit{Source reference.} Original lemma 02LK: sources/stacks-more-morphisms/pinned-original-source.tex, source lines 11924--11924. Navigation pointer only: the counted proof cites this statement by \textbackslash{}ref as a standing input; the deduction the counted proof has to make is carried out inside the counted span, and no part of that deduction is delegated to this pointer.\par\medskip
\ifdefined\hypertarget\hypertarget{sourcepointer-3}{}\fi
\noindent\textit{Source reference.} Original lemma 0F3G: sources/stacks-more-morphisms/pinned-original-source.tex, source lines 22514--22514. Navigation pointer only: the counted proof cites this statement by \textbackslash{}ref as a standing input; the deduction the counted proof has to make is carried out inside the counted span, and no part of that deduction is delegated to this pointer.\par\medskip
\ifdefined\hypertarget\hypertarget{sourcepointer-4}{}\fi
\noindent\textit{Source reference.} Original lemma lemma-weighting-specialization: sources/stacks-more-morphisms/pinned-original-source.tex, source lines 22469--22469. Navigation pointer only: the counted proof cites this statement by \textbackslash{}ref as a standing input; the deduction the counted proof has to make is carried out inside the counted span, and no part of that deduction is delegated to this pointer.\par\medskip
\ifdefined\hypertarget\hypertarget{sourcepointer-5}{}\fi
\noindent\textit{Source reference.} Original reference 03HX: focused-seeds/geometry/S9-morph-weighting-dense-set/ref-tags.txt, source lines 5--5. Navigation pointer only: the counted proof cites this statement by \textbackslash{}ref as a standing input; the deduction the counted proof has to make is carried out inside the counted span, and no part of that deduction is delegated to this pointer.\par\medskip
\ifdefined\hypertarget\hypertarget{sourcepointer-6}{}\fi
\noindent\textit{Source reference.} Original reference topology-lemma-dense-in-constructible: focused-seeds/geometry/S9-morph-weighting-dense-set/ref-tags.txt, source lines 6--6. Navigation pointer only: the counted proof cites this statement by \textbackslash{}ref as a standing input; the deduction the counted proof has to make is carried out inside the counted span, and no part of that deduction is delegated to this pointer.\par\medskip
\begin{lemma}
\label{lemma-weighting-on-dense-set}
Let $f : X \to Y$ be a locally quasi-finite morphism of schemes.
Let $w : X \to \mathbf{Z}$ be a weighting of $f$.
Let $E \subset X$ be a dense subset.
Assume $w(x) = 0$ for all $x \in E$. Then $w = 0$.
\end{lemma}

\begin{proof}
We suggest reading Lemma \ref{lemma-weighting-specialization} first.
Let $x \in X$ with image $y \in Y$. We will show that $w(x) = 0$. By
Lemma \ref{lemma-etale-makes-quasi-finite-finite-at-point}
we can choose a diagram
$$
\xymatrix{
U \ar[r] \ar[dr]_\pi & X_V \ar[d] \ar[r] & X \ar[d]^f \\
& V \ar[r]^\varphi & Y
}
$$
with $V \to Y$ \'etale, with a point $v \in V$ mapping to $y$,
with $U \subset Y_V$ open, such that $U \to V$ is finite, and
such that there exists an unique point $u \in U$ mapping to $v$
which moreover maps to $x$ in $X$. We may and do assume that $V$
(and hence $U$) is affine. By Definition \ref{definition-weighting}
the function $\int_\pi (w \circ h)$ is locally constant where
$h : U \to X$ is the composition of the top horizontal arrows.
The set of points $U_0$ where $w \circ h$ is $0$ is a
constructible subset of $U$ by Lemma \ref{lemma-jumps-w}
and dense in $U$ because it contains $h^{-1}(E)$ (use that $h$ is open).
By Topology, Lemma \ref{topology-lemma-dense-in-constructible}
we see that $U_0$ contains a dense open,
i.e., the set of points where $w \circ h$ is nonzero is contained
in a nowhere dense closed subset $Z \subset U$. Then $\pi(Z)$
is a nowhere dense closed subset of $V$ by
Morphisms, Lemma \ref{morphisms-lemma-image-nowhere-dense-finite}.
Hence $\int_\pi (w \circ h)$ is zero on a dense open hence
vanishes everywhere.
However, the value at $v$ is $w(x)[\kappa(u):\kappa(v)]_s$
and we conclude.
\end{proof}

\end{document}
